% Generated draft: requires mathematical and source review.
\documentclass[11pt]{article}
\usepackage[margin=25mm]{geometry}
\usepackage{fontspec}
\setmainfont{Latin Modern Roman}
\usepackage{amsmath,amssymb,mathtools,unicode-math}
\setmathfont{Latin Modern Math}
\usepackage{enumitem,needspace,xurl,hyperref}
\setlength{\parindent}{0pt}
\setlength{\parskip}{5pt}
\setlength{\emergencystretch}{3em}
\newcommand{\N}{\mathbb{N}}
\newcommand{\Z}{\mathbb{Z}}
\newcommand{\Q}{\mathbb{Q}}
\newcommand{\R}{\mathbb{R}}
\newcommand{\C}{\mathbb{C}}
\newcommand{\F}{\mathbb{F}}
\newcommand{\A}{\mathbb{A}}
\newcommand{\PP}{\mathbb{P}}
\newcommand{\E}{\mathbb{E}}
\newcommand{\eps}{\varepsilon}
\newcommand{\dd}{\,\mathrm{d}}
\begin{document}
\section*{UnsolvedMath ID: 11100064}
\textbf{Source problem code:} catalogue entry
\subsection{Definitions and mathematical statement}
For a spectrum $E$, its Bousfield class $\langle E\rangle$ is determined by its acyclics: $\langle E\rangle=\langle F\rangle$ means $E\wedge Y\simeq0$ if and only if $F\wedge Y\simeq0$ for every spectrum $Y$. To kill $X$ means $E\wedge X\simeq0$.

Hovey's Miscellaneous item 2 first asks to classify the Bousfield classes realized by $E_\infty$ ring spectra. It then proposes the following explicitly $p$-local conjecture: if an $E_\infty$ ring spectrum $E$ kills a nontrivial finite spectrum, then $\langle E\rangle=\langle E(n)\rangle$ for some $n$, where $E(n)$ denotes Johnson--Wilson theory. The source does not give the range of $n$; in the conventional reading $n=0$ allows rational theory. It also does not explicitly exclude the zero ring spectrum, which would be an immediate exception unless a convention accounts for the zero Bousfield class. Thus a meaningful nonzero formulation is conditional on that missing qualification:
\[
 E\not\simeq0,\quad\exists X\not\simeq0\text{ finite $p$-local},\ E\wedge X\simeq0
 \quad\Longrightarrow\quad
 \exists n\geq0:\langle E\rangle=\langle E(n)\rangle.
\]
No numerical height bound in terms of $X$ is supplied.

The source warns that dualizing the diagonal of a suspension spectrum must include the disjoint basepoint. For nonempty finite $X$, the dual $D\Sigma^\infty X_+$ has a sphere retract and hence the sphere Bousfield class; the empty-space and infinite-space dual conventions need separate treatment. This prevents the suggested example from automatically producing other classes.

Hovey recalls a tentative attribution to Jim McClure that making a quotient such as $MU/p$ into an $E_\infty$ ring may require also killing higher $v_i$. This is hearsay without a specified quotient or generators, not an established universal relation. He cites Goerss--Hopkins for Morava-$E$ $E_\infty$ structures and says these realize the $E(n)$ classes. Morava-$E$ spectra and Morava $K(n)$ Bousfield classes must not be silently identified. The registry's vague modern rigidity/telescope discussion is truncated and supplies no exact theorem resolving the conjecture.
\subsection{Short English statement}
Classify Bousfield classes of $E_\infty$ rings, especially the proposed $p$-local statement that a nonzero one with a nontrivial finite acyclic has Johnson--Wilson Bousfield class. Zero-ring and height conventions remain unresolved.
\subsection{Sources}
\begin{itemize}
\item Local registry: \url{data/unsolvedmath/problems.json}, permanent UnsolvedMath ID $11100064$, exact source code \texttt{catalogue entry}, attributed to Mark Hovey.
\item \emph{Hovey's Algebraic Topology Problem List}, Miscellaneous problems item 2: \url{https://www-users.cse.umn.edu/~tlawson/hovey/}. Supplied access date: 2026-07-29; extraction: html-ordered-list; 65 top-level items in nine category pages.
\item Source attributions: Goerss--Hopkins for Morava-$E$ structures, and a tentative recollection ``Jim McClure?'' for the $MU/p$ observation. No precise articles are supplied.
\item The registry mentions Barthel--Heard--Sanders and unspecified telescope, rigidity and module results without exact references or hypotheses; its final summary is truncated.
\end{itemize}
\end{document}
